\documentclass[11pt,a4paper]{article}

\usepackage[utf8]{inputenc}
\usepackage[T1]{fontenc}
\usepackage[backend=biber,style=numeric,sorting=none]{biblatex}
\usepackage[margin=2.5cm]{geometry}
\usepackage[hidelinks]{hyperref}
\usepackage{amsmath}
\usepackage{amssymb}
\usepackage{graphicx}

\addbibresource{references.bib}

\title{Build-System Hermeticity as a First-Class Concern}
\author{Jane Engineer \\ \texttt{jane@example.com}}
\date{\today}

\begin{document}
\maketitle

\begin{abstract}
We examine the cost of build-system non-hermeticity in industrial
software projects, surveying recent practice from
the literature~\cite{morgenthaler2012searching, esfahani2016cloudbuild}
and contrasting it with content-addressed approaches popularised by
modern build systems~\cite{erickson2018bazel}. We conclude that
toolchain-level hermeticity, while occasionally painful to adopt, is
a forcing function for several adjacent quality properties.
\end{abstract}

\section{Introduction}
\label{sec:intro}

In a 2012 survey of build practice at Google, Morgenthaler et al.
report that an estimated 36\% of engineering time was lost to
build-related friction~\cite{morgenthaler2012searching}. A
subsequent generation of build systems --- including Bazel,
Buck, and Pants --- has aimed to reduce that friction in part by
treating hermeticity as a first-class concern.

\section{Background}
\label{sec:background}

Hermeticity in this paper means: \emph{given identical inputs, a
build action produces identical outputs, regardless of when, where,
or by whom it is run}. Content-addressed inputs and sandboxed
execution are the two main mechanisms by which modern build systems
realise this property~\cite{erickson2018bazel}.

The claim in
Section~\ref{sec:cost} is that hermeticity also implies several
adjacent properties (cacheability, distributed execution,
reproducibility) that compound to large productivity gains.

\section{The Cost and Benefit}
\label{sec:cost}

\dots

\printbibliography

\end{document}
